\documentclass[paper=a4paper,fontsize=10.5pt]{jlreq}
\usepackage{booktabs}
\usepackage{amsmath}
\usepackage{graphicx}
\usepackage[hidelinks]{hyperref}

\title{ECAPA-TDNNのモバイル向け軽量化と\\オンデバイス話者照合PoC}
\date{2026年11月3日}

\begin{document}
\maketitle

\begin{abstract}
SpeechBrainの学習済みECAPA-TDNNを，前処理を含む単一のONNXモデルへ変換し，ONNX Runtimeの静的INT8量子化によって軽量化した．
モデルサイズは84.1\,MBから25.8\,MBに減少し，CPU上の推論時間は約5分の1となった一方，LibriSpeech上の等誤り率（EER）に劣化は見られなかった．
さらに，このモデルを用いて登録音声と発話音声の同一性を判定するFlutterアプリを試作した．
\end{abstract}

\section{目的}
話者照合モデルECAPA-TDNNをスマートフォン上で動作させ，事前に登録した音声と新たな発話が同一人物によるものかをオンデバイスで判定することを目的とする．
アプリはFlutterで実装し，推論にはONNX Runtimeを用いる．

\section{手法}

\subsection{ベースモデル}
SpeechBrainが公開する\texttt{spkrec-ecapa-voxceleb}（VoxCelebで学習，パラメータ数約2,080万，埋め込み192次元）を用いた．
本報告ではモデル構造の変更や再学習は行っていない．

\subsection{前処理を同梱したONNX化}
元のモデルは80次元の対数メルフィルタバンク（Fbank）を入力とする．
アプリ側で同じ特徴量抽出を再実装すると，実装差による精度劣化の原因となるため，Fbankをモデルに含めた．
具体的には，ハミング窓をかけたDFT基底を重みとする1次元畳み込みでSTFTを表現し，パワースペクトルへのメル行列の乗算，対数化，ダイナミックレンジ80\,dBでのクリップ，発話単位の平均除去をすべて基本演算で記述した．
STFT演算子を用いないため，古いランタイムや各種実行プロバイダでも動作しやすい．
得られたモデルの入出力は次のとおりである．
\begin{itemize}
  \item 入力：16\,kHzモノラル波形（float32，長さ可変）
  \item 出力：L2正規化済み192次元埋め込み
\end{itemize}
SpeechBrain本来の推論結果との一致を確認したところ，コサイン類似度は1.000000であった．
また，1.5秒，3秒，7秒の入力でPyTorchとONNX Runtimeの出力差は最大$1.9\times10^{-7}$であり，可変長入力が正しく扱えることを確認した．

\subsection{量子化}
ONNX Runtimeの量子化ツールを用いて次の2方式を比較した．
\begin{itemize}
  \item \textbf{動的量子化}：重みのみをINT8化し，活性のスケールは実行時に求める．校正データは不要である．
  \item \textbf{静的量子化（QDQ形式）}：重みをチャネル単位でINT8化し，活性は校正データから求めた範囲（99.99パーセンタイル）でUINT8化する．
\end{itemize}
いずれも前処理部分（DFT畳み込み，対数メル）は精度保持のためFP32のまま残した．

\subsection{照合方式}
登録時に3回の発話からそれぞれ埋め込みを求め，その平均を再正規化して登録声紋とする．
照合時は発話の埋め込みと登録声紋のコサイン類似度$s$を求め，$s \geq \theta$ならば本人と判定する．
入力前にはフレームエネルギーに基づく簡易な無音除去を行い，有声区間が1.5秒未満の発話は棄却する．

\section{実験}

\subsection{条件}
評価にはLibriSpeechを用いた．
test-cleanから40話者$\times$12発話（3秒以上）を抽出したものに，動作確認用に取得していたdev-cleanの1話者73発話が加わった計41話者553発話を用意した．
このうち10話者を校正専用とし（各2発話を使用），残る31話者433発話を評価に用いた．
各発話の先頭4秒から埋め込みを求め，全組合せ93,528試行（本人4,608，他人88,920）でEERを算出した．
推論時間はx86 CPU上で1スレッド，3秒入力に対する10回の平均である．

\subsection{結果}
結果を表\ref{tab:result}に示す．

\begin{table}[htbp]
  \centering
  \caption{量子化方式ごとの比較}
  \label{tab:result}
  \begin{tabular}{lrrrrr}
    \toprule
    モデル & サイズ [MB] & EER [\%] & EER閾値 & FP32とのcos & 推論 [ms] \\
    \midrule
    FP32            & 84.1 & 0.80 & 0.327 & 1.0000 & 225.6 \\
    INT8 動的       & 22.0 & 0.78 & 0.334 & 0.9947 & 320.1 \\
    INT8 静的（QDQ）& 25.8 & 0.78 & 0.329 & 0.9957 & 45.6 \\
    \bottomrule
  \end{tabular}
\end{table}

静的量子化はサイズを約3分の1に，推論時間を約5分の1にしつつ，EERを維持した．
動的量子化はサイズでは最小であるが，畳み込み主体のモデルでは実行時の量子化処理が負担となり，FP32より遅くなった．
推論時間には測定ごとのばらつきがあり，別の測定ではそれぞれ393.7，607.8，77.8\,msであったが，相対的な傾向は変わらない．
以上より，アプリには静的量子化モデルを採用した．

なお，Flutterパッケージに同梱されるONNX Runtime 1.15.1でも，FP32モデルと静的量子化モデルが読み込めて，正規化済み埋め込みが得られることを確認した．

\subsection{実装上の問題と対処}
\begin{enumerate}
  \item \textbf{可変長入力での破綻}：全演算子を量子化対象にすると，Attentive Statistics Poolingのマスク長計算まで量子化され，校正時と異なる長さの入力でエラーとなった．量子化対象をConvとMatMulに限定して解決した．
  \item \textbf{ランタイムのクラッシュ}：ONNX Runtime 1.31のQDQ融合最適化が，最終全結合層（Padを経由する1×1畳み込み）でクラッシュした．当該ノード（約120万パラメータ）のみFP32で残して回避した．
  \item \textbf{旧ランタイムでの読込失敗}：量子化ツールが付与する未使用のopset宣言（\texttt{ai.onnx.ml} v5など）により，ONNX Runtime 1.15がロードを拒否した．未使用の宣言を削除して解決した．
\end{enumerate}

\section{Flutterアプリ}
\texttt{onnxruntime}パッケージ（1.4.1）と\texttt{record}パッケージを用い，Android／iOS向けに実装した．
マイクから16\,kHzのPCMを取得してfloatに変換し，無音除去の後にモデルへ入力する．
登録声紋は端末内に保存し，判定閾値は画面上のスライダーで調整できる．
OSのノイズ抑圧や自動利得制御は声質を変化させるため無効にしている．
静的解析とユニットテストは通過しているが，APKのビルドおよび実機での動作確認は未実施である．

\section{課題と今後}
\begin{itemize}
  \item \textbf{評価の妥当性}：評価は英語の読み上げ音声によるもので，規模も小さい．日本語，実機のマイク，雑音環境での評価と閾値の再調整が必要である．
  \item \textbf{実機性能}：スマートフォン上での推論時間，メモリ使用量，ランタイムを含むアプリサイズを測定する必要がある．
  \item \textbf{さらなる軽量化}：チャネル数を1024から512に減らした構造を，再学習または知識蒸留で学習すれば，6〜8\,MB程度への縮小が見込める．WeSpeakerなどの小型学習済みモデルとの比較も考えられる．
  \item \textbf{安全性}：録音音声の再生によるなりすましへの対策が別途必要である．
\end{itemize}

\section{まとめ}
ECAPA-TDNNを前処理込みの単一ONNXモデルとし，静的INT8量子化を施すことで，精度を保ったまま大幅な軽量化と高速化を達成した．
また，これを用いたFlutterによるオンデバイス話者照合の試作を行った．
今後は実機と実環境の音声による評価を進める．

\end{document}
